\section{Conclusion and Future Horizons}
\label{sec:conclusion}

In this work, we introduced \textbf{ThermoMerge} (Thermodynamic Model Merging), a paradigm-shifting framework that rethinks model merging through the lens of statistical physics, statistical mechanics, and thermodynamics. By moving away from the traditional, rigid, and cold assumption of static linear parameter interpolation in Euclidean space, we modeled the merging process as a dynamic, finite-temperature thermal-equilibrium adaptation. We showed that mapping classification logits to negative energies in a canonical Boltzmann ensemble provides a mathematically rigorous foundation for model merging.

Through our Thermodynamic Annealing Schedule (TAS), we demonstrated that simulated thermal cooling flattens highly non-convex multi-task loss frontiers, permitting trainable merging parameters to escape frustrated local minima. Furthermore, we proved that minimizing the Helmholtz Free Energy Discrepancy (F-Min) acts as a natural physical regularizer that is mathematically proportional to the temperature-scaled Kullback-Leibler divergence. Empirically, ThermoMerge mitigates the severe transductive overfitting pathology of the \emph{Overfitting-Optimizer Paradox} on simple visual tasks, achieving highly competitive multi-task accuracies compared to unregularized adaptive methods, and establishing a stable, physically grounded framework for multi-task adaptation.

\subsection{Future Horizons}
As **The Visionary**, we believe that the successful marriage of thermodynamics and model merging is not merely a specialized solution for parameter fusion, but the first step toward a broader, physical paradigm of general intelligence. We outline several highly promising, paradigm-shifting research directions:
\begin{itemize}
    \item \textbf{Dynamic Thermal Network Adaptation and Weight-Space Heat Capacities:} Extending local task-wise thermal coupling to individual layers, or exploring the application of layer-specific temperatures directly to weight updates (e.g., scaling parameter updates by layer-wise physical heat capacities). This can provide crucial localized regularization, especially on complex multi-channel color tasks where representations are fragile and prone to interference.
    \item \textbf{Scaling to Modern Foundation Models:} Verifying and scaling these thermodynamic principles on pre-trained foundation models (e.g., CLIP ViT-B/32 or ResNet-50) across standard multi-task benchmarks to demonstrate that the physics-grounded optimization generalizes to massive, overparameterized transformer architectures. We provide a complete, mathematically concrete roadmap detailing PEFT/adapter parameterization, multimodal logit-to-energy formulations, prediction caching, and layer-wise heat capacities in Appendix~\ref{sec:scaling_roadmap}.
    \item \textbf{Ecological and Symbiotic Multi-Agent Merging:} Treating multiple foundation models as distinct biological species in a shared neural ecosystem, where model fusion is governed by symbiotic and evolutionary thermodynamics rather than manual parameter routing.
    \item \textbf{Quantum Ensemble Merging:} Elevating our classical thermodynamic framework to a quantum statistical paradigm, modeling parameters as wavefunctions and tasks as unitary transformation operators, enabling quantum-like superposition and collapse of multi-task weights.
\end{itemize}

By bridging deep learning with the fundamental laws of nature, ThermoMerge paves a bold, new pathway toward fluid, self-crystallizing, and physically grounded artificial generalists.
